\chapter{문제 정식화}
\label{ch:problem}

\section{텔레메트리의 형식적 모델}
\label{sec:problem-telemetry}

한 경기 $m$의 timeline을 다음 두 열로 모형화한다.

\begin{itemize}[leftmargin=1.5em]
  \item \textbf{프레임 열} $\mathcal{F}_m = \{F_0, F_1, \dots, F_{T}\}$. 프레임 $F_k$는 시각 $t_k = 60{,}000\,k$ ms에서의 관측이며, 참가자 $i \in \{1,\dots,10\}$에 대한 상태 벡터 $s_i(t_k)$(위치 $\mathrm{pos}_i(t_k) \in [0,16000]^2$ 포함)와 팀 단위 집계를 담는다.
  \item \textbf{사건 열} $\mathcal{E}_m = \{u_1, u_2, \dots\}$. 사건 $u$는 종류 $\mathrm{type}(u)$, 타임스탬프 $t(u) \in \mathbb{N}$ (ms), 팀 $\mathrm{team}(u) \in \{100, 200\}$, 그리고 종류에 따라 위치 $\mathrm{pos}(u)$와 참가자 집합 $\mathrm{part}(u)$를 갖는다.
\end{itemize}

킬 사건 집합을 $\mathcal{K}_m = \{u \in \mathcal{E}_m : \mathrm{type}(u) = \texttt{CHAMPION\_KILL}\}$이라 하고 시각 순으로 정렬한다.

\section{교전의 정의}
\label{sec:problem-engagement}

\begin{definition}[킬 조건부 교전]
\label{def:engagement}
경기 $m$의 교전은 튜플 $g = (t_e, t_{k_1}, t_{k_L}, c, \mathcal{K}_g)$이다. 여기서 $\mathcal{K}_g \subseteq \mathcal{K}_m$은 서로 시간적으로 연속(간격 $\le \delta_{\mathrm{gap}}$)이고 공간 지름이 $D_{\max}$ 이하인 킬의 집합, $t_{k_1}, t_{k_L}$은 그 첫·마지막 킬 시각, $c = \mathrm{pos}(k_1)$은 첫 킬 위치(교전 중심), $t_e = t_{k_1} - \delta_{\mathrm{pre}}$는 \emph{온셋}(engage time)이다. 교전은 온셋 시각에 양 팀 각각 생존 챔피언 두 명 이상이 중심 $c$로부터 반경 $r_{\mathrm{val}}$ 안에 있을 때에만 유효하다.
\end{definition}

이 정의는 Schubert 등 \cite{schubert2016esports}의 encounter 개념(공간적 근접과 전투 사건에 의한 집단 상호작용)을 분 단위 텔레메트리에 맞게 재구성한 것이다. 리플레이 기반 정의가 피해 교환을 직접 관측하는 것과 달리, 본 정의는 관측 가능한 유일한 전투 신호인 \emph{킬}을 씨앗으로 삼는다. 따라서 킬이 발생하지 않은 대치(standoff)는 정의상 교전에 포함되지 않으며, 이는 본 논문의 명시적 범위 제한이다.

\section{예측 문제}
\label{sec:problem-prediction}

교전 $g$가 주어졌을 때, \textbf{관측 창} $W_{\mathrm{obs}} = [t_e - c_{\mathrm{ctx}},\, t_e)$와 \textbf{라벨 창} $W_{\mathrm{lab}} = [t_e,\, t_e + h)$를 정의한다. 본 논문은 $c_{\mathrm{ctx}} = h = 30{,}000$ ms를 사용하며, 관측 창을 길이 $b = 5{,}000$ ms의 $L = 6$개 구간으로 나눈다.

\begin{definition}[교전 결과 예측 문제]
관측 창에서 얻은 다중 모달 표현
\begin{equation}
  X_g = \big(X^{\mathrm{node}} \in \R^{L\times N\times F_{\mathrm{node}}},\; X^{\mathrm{glob}} \in \R^{L\times F_{\mathrm{glob}}},\; X^{\mathrm{ev}} \in \R^{L\times F_{\mathrm{ev}}},\; E \in \R^{K\times D_e}\big)
\end{equation}
로부터, 라벨 창의 사건만으로 결정되는 이진 라벨 $y_g = \Lambda(\mathcal{E}_m \cap W_{\mathrm{lab}}) \in \{0,1\}$ ($1$: 블루 팀 승)을 예측하는 점수 함수 $f(X_g) \approx P(y_g = 1 \mid X_g)$를 학습한다.
\end{definition}

$\Lambda$는 제 \ref{ch:label} 장의 교환가치 라벨 함수이고, $F_{\mathrm{node}} = 76$, $F_{\mathrm{glob}} = 26$, $F_{\mathrm{ev}} = 44$, $K = 64$, $D_e = 12$이다.

\section{무누출 계약}
\label{sec:problem-contract}

시간 해상도 불일치(60 초 프레임 대 밀리초 사건) 때문에 관측 창의 상태 값을 어떤 프레임에서 가져올지가 곧 누출 여부를 결정한다. 본 논문은 다음 세 계약을 모든 모델 입력에 강제하며, 각 계약은 단위 시험으로 검증된다.

\begin{assumption}[계단 유지 계약]
\label{asm:stephold}
구간 $\ell$의 중점 $q_\ell = t_e - c_{\mathrm{ctx}} + (\ell + \tfrac12) b$에서의 노드·전역 상태는 $q_\ell$ \emph{이전}의 마지막 프레임 $F_{k^*}$, $k^* = \max\{k : t_k < q_\ell\}$의 값을 그대로 사용한다(piecewise-constant / forward-fill). 프레임 사이를 보간하지 않는다.
\end{assumption}

\begin{assumption}[온셋 상한 계약]
\label{asm:onsetcap}
어떤 구간에서도 $t_k \ge t_e$인 프레임은 사용하지 않는다. 즉 $k^* \le \max\{k : t_k < t_e\}$이다.
\end{assumption}

\begin{assumption}[킬 좌표 배제 계약]
\label{asm:nokillpos}
교전 국소화 단계에서만 사용되는 5 초 위치 격자(제 \ref{sec:loc-grid} 절)와 킬 좌표는 모델 입력에 포함되지 않는다. 사건 특징은 관측 창 $[b_0, b_1)$ 안에서 발생한 사건의 계수만을 사용한다.
\end{assumption}

가정 \ref{asm:stephold}는 아이템 구매·버프 만료처럼 상태 전이가 이산적인 게임 상태에 대해 선형 보간이 존재하지 않는 매끄러움을 만들어내는 것을 막고, 가정 \ref{asm:onsetcap}은 교전이 이미 진행 중인 프레임의 체력·생존 정보가 입력으로 새는 것을 막는다. 관측 창의 길이가 30 초이고 프레임 간격이 60 초이므로, 대부분의 교전에서 여섯 구간은 동일한 프레임 하나(또는 둘)를 참조한다. 이것이 본 논문이 ``분 단위'' 제약을 강조하는 이유이며, 관측 창 안의 시간 변화는 사실상 사건 계수와 (드물게) 프레임 경계에서만 발생한다.

\section{표기 요약}
\label{sec:problem-notation}

\begin{table}[H]
  \centering
  \caption{주요 기호}
  \label{tab:notation}
  \small
  \begin{tabular}{@{}lll@{}}
    \toprule
    기호 & 의미 & 값 \\
    \midrule
    $N$ & 참가자 수 & 10 \\
    $L$ & 관측 창의 구간 수 & 6 \\
    $b$ & 구간 길이 & 5{,}000 ms \\
    $c_{\mathrm{ctx}}$ & 관측 창 길이 & 30{,}000 ms \\
    $h$ & 라벨 창(예측 지평) 길이 & 30{,}000 ms \\
    $\delta_{\mathrm{gap}}$ & 킬 시간 군집 간격 & 18{,}000 ms \\
    $D_{\max}$ & 킬 군집의 최대 공간 지름 & 4{,}000 단위 \\
    $\delta_{\mathrm{pre}}$ & 첫 킬 이전 온셋 오프셋 & 10{,}000 ms \\
    $r_{\mathrm{val}}$ & 온셋 검증 반경 & 1{,}800 단위 \\
    $n_{\min}$ & 팀당 최소 생존 인원 & 2 \\
    $\delta_{\mathrm{merge}}, r_{\mathrm{merge}}$ & 인접 후보 병합 간격·반경 & 15{,}000 ms, 2{,}000 단위 \\
    $T_{\max}$ & 병합 후 최대 교전 길이 & 60{,}000 ms \\
    $F_{\mathrm{node}}, F_{\mathrm{glob}}, F_{\mathrm{ev}}$ & 노드·전역·사건 특징 차원 & 76, 26, 44 \\
    $K, D_e$ & 사건 토큰 수·연속 특징 차원 & 64, 12 \\
    $\beta$ & 라벨 주의 온도 & 2.0 \\
    \bottomrule
  \end{tabular}
\end{table}
